\subsection{The exact injectivity threshold for the explicit threefold}

\begin{lemma}[The corrected dual symmetric-power calculation]
\label{lem:OC-threefold-negative-degrees}
\status{Proved; independent root and reduction-descent proof, primary-source, and computation audits completed.}
For the fixed $S,L_2,R,L,X,A$ of
Proposition~\ref{prop:OC-Tango-threefold} and
Lemma~\ref{lem:OC-threefold-multiroot},
\[
 h^1(R,L^{-1})=360,\qquad h^1(R,L^{-2})=1,\qquad
 H^1(R,L^{-r})=0\quad(r\geq3).
\]
Consequently
\[
 h^1(X,A^{-1})=1203,\qquad h^1(X,A^{-2})=1,\qquad
 H^1(X,A^{-d})=0\quad(d\geq3).
\]
\end{lemma}
\begin{proof}
Write $B=L_2$ and let $\alpha\in H^1(S,B^{-10})$ be the nonsplit Tango extension class defining $V=V_2$.  A useful preliminary fact is
\begin{equation}\label{eq:OC-Tango-nonannihilation}
 H^0(S,V^\vee\otimes B^e)=0\qquad(0\leq e<10).
\end{equation}
The extension sequence identifies this group with the kernel of
\[
 H^0(S,B^e)\longrightarrow H^1(S,B^{e-10}),\qquad s\longmapsto s\alpha,
\]
since $H^0(B^{e-10})=0$.  Under the exact-differential description of the Frobenius kernel, multiplication by $s$ replaces the nonzero rational differential $\eta$ representing $\alpha$ by $s^2\eta$.  This is nonzero whenever $s\ne0$.  Moreover the connecting map is injective for the ample effective divisor $(10-e)L_2$: the preceding group $H^0(B^{-2(10-e)})$ is zero.  Thus $s\alpha\ne0$ for every nonzero $s$, proving \eqref{eq:OC-Tango-nonannihilation}.  This is the naturality of the exact sequence in \cite[Theorem~1 and Lemma~1.4, pp.~1, 4--5]{Mukai2011Kodaira}; no assertion about multiplication on arbitrary cohomology is being made.

For $r\geq1$, put
\[
 m_j(r)=\left\lceil\frac{2(r+j)}5\right\rceil,
 \qquad e_j(r)=4j-r-10,\qquad0\leq j\leq4.
\]
All eigensheaves of $\psi_{2*}L^{-r}$ have strictly negative relative degree.  Thus $g_*L^{-r}=0$, and relative duality with Leray gives
\begin{equation}\label{eq:OC-correct-negative-R1}
 H^1(R,L^{-r})=
 \bigoplus_{m_j(r)\geq2}
 H^0\!\left(S,\bigl(\operatorname{Sym}^{m_j(r)-2}V\bigr)^\vee
                                      \otimes B^{e_j(r)}\right).
\end{equation}
The dual is taken \emph{after} the symmetric power.  In characteristic two it must not be replaced by $\operatorname{Sym}^{m_j(r)-2}(V^\vee)$.

For $r\geq7$, all $e_j(r)\leq6-r<0$.  The dual symmetric power has a filtration with quotients $B^{-10a}$, so every summand in \eqref{eq:OC-correct-negative-R1} vanishes.
At $r=3$, only $j=4$ can contribute, and its group is $H^0(V^\vee\otimes B^3)=0$ by \eqref{eq:OC-Tango-nonannihilation}.
At $r=4,5,6$, the only possible contribution is
\[
              H^0\!\left(S,(\operatorname{Sym}^2V)^\vee
                                                   \otimes B^{6-r}\right).
\]
The Frobenius splitting implies the genuine decomposition
\begin{equation}\label{eq:OC-symmetric-square-splitting}
                         \operatorname{Sym}^2V\simeq V\oplus B^{20}.
\end{equation}
To check it, write the global subline generator as $u$ and local complementary vectors as $v_i$, with $v_j=a_{ij}v_i+c_{ij}u$.  The subbundle $uV$ of $\operatorname{Sym}^2V$ has basis $u^2,uv_i$ and is isomorphic to $V$.  A splitting of $F^*V$ supplies local combinations $v_i^2+h_i u^2$ gluing as the line $B^{20}$; they complement $uV$.  Equivalently, the transition matrix in this basis is the direct sum of the original extension matrix and $a_{ij}^2$.  Dualizing \eqref{eq:OC-symmetric-square-splitting} and using \eqref{eq:OC-Tango-nonannihilation} proves vanishing at $r=4,5,6$.

At $r=2$, the $j=3$ term is $H^0(\OO_S)=k$, the $j=4$ term is $H^0(V^\vee\otimes B^4)=0$, and the other terms vanish by negativity.  At $r=1$, the only terms are $H^0(B)$ and $H^0(B^5)$.
The first normalized cover gives
\[
 f_*B=N_0\oplus N_0^{21},\qquad
 f_*B^5=\bigoplus_{i=0}^{5}N_0^{5+20i}.
\]
For $0\leq b\leq105$, $h^0(C,N_0^b)=b+1$: a function with its only pole at infinity lies in
$k[x,y]/(y^2+y-x^{843})$ and is uniquely $a(x)+y b(x)$; its two summands have even and odd pole orders, and a nonzero second summand has pole order at least $843$.  All pole bounds here are at most $210$.  Therefore
\[
 h^0(B)=2+22=24,\qquad
 h^0(B^5)=\sum_{i=0}^{5}(6+20i)=336,
\]
and $h^1(R,L^{-1})=360$.

The multiroot eigensheaves now give
\[
 H^1(X,A^{-d})=\bigoplus_{0\leq a_i<445}
                  H^1(R,L^{-(d+\sum a_i)}).
\]
Only exponents one and two contribute.  There are $843$ tuples of sum one, so the dimensions are $360+843=1203$ at $d=1$, one at $d=2$, and zero thereafter.
\end{proof}



\begin{lemma}[All one-forms with the required negative twists]
\label{lem:OC-threefold-differential-degrees}
\status{Proved; independent root and reduction-descent proof, primary-source, and computation audits completed.}
Let $\vartheta:L^4\to\Omega_R^1$ be the nowhere-zero Tango map from the second iteration, and also write $\vartheta:A^4\to\Omega_X^1$ for its pullback followed by the differential map.  For every integer $r\geq2$ and $d\geq2$,
\[
 H^0(R,\Omega_R^1\otimes L^{-r})
          =\vartheta\,H^0(R,L^{4-r}),
 \qquad
 H^0(X,\Omega_X^1\otimes A^{-d})
          =\vartheta\,H^0(X,A^{4-d}).
\]
In particular both sides vanish when the exponent of the negative twist is at least five.  The map $\vartheta$ remains nonzero after restriction to any smooth ample hypersurface $Y\in|A^d|$.
\end{lemma}
\begin{proof}
Write again $B=L_2$ and $V=V_2$.  The first Tango step supplies the subbundle $B^{20}\hookrightarrow\Omega_S^1$, with locally free quotient $Q$ of rank one.  On $P_2=\mathbf P_S(V)$, the logarithmic cotangent bundle for the disjoint smooth branches $E_0,G$ has a filtration with successive quotients
\begin{equation}\label{eq:OC-log-cotangent-filtration}
 \OO_{P_2}(2),\qquad\pi^*Q,\qquad
 \Omega^1_{P_2/S}(E_0)=\OO_{P_2}(-1)\otimes B^{10}.
\end{equation}
Indeed $G$ locally has equation $S^2-f_i=0$, where $df_i$ spans the distinguished base line.  The only new horizontal logarithmic generator is $df_i/(S^2-f_i)$, so the horizontal kernel is an extension of $\pi^*Q$ by
$\pi^*B^{20}(G)=\OO_{P_2}(2)$.
The relative quotient acquires exactly a simple pole along $E_0$ and none along $G$, since $d_{P_2/S}(S^2-f_i)=0$.  At infinity the regular equation of $G$ is $1-f_i z^2$, and the same horizontal generator is $df_i/(1-f_i z^2)$; the relative logarithmic generator is $dz/z$.  These descriptions verify the filtration on both charts and their overlaps.

The normalized tame cover identifies the pulled-back logarithmic cotangent bundle with the logarithmic cotangent bundle upstairs, and hence
\[
             \Omega_R^1\hookrightarrow
             \psi_2^*\Omega_{P_2}^1(\log(E_0+G)).
\]
Use the negative eigensheaves and the integers $m_j(r)$ from the preceding lemma.  The last two graded pieces after twisting have relative degrees $-m_j(r)$ and $-1-m_j(r)$, so they have no sections.  The first piece is
\begin{equation}\label{eq:OC-first-log-piece}
                 \OO_{P_2}(2-m_j(r))\otimes B^{4j-r}.
\end{equation}
For $r\geq6$ all these relative degrees are negative.  At $r=5$ only $j=0$ has degree zero, and its base twist is $B^{-5}$; again there are no sections.

At $r=4$, the sole possible nonzero term is $j=1$, giving $H^0(\OO_S)$.  At $r=3$, the possible terms give $H^0(B)\oplus H^0(B^5)=H^0(R,L)$.  At $r=2$, they give
\[
 H^0(S,V\otimes B^{-2})\ \oplus\ H^0(B^2)
                      \oplus H^0(B^6)\oplus H^0(B^{10}).
\]
The first group is zero.  Indeed its extension sequence has connecting map
$H^0(B^8)\to H^1(B^{-2})$, $s\mapsto s\alpha$, which is injective by the nonannihilation argument \eqref{eq:OC-Tango-nonannihilation} at $e=8$.
The other three groups equal $H^0(R,L^2)$ by the positive normalized eigensheaves at $a=4$.
These calculations bound $h^0(\Omega_R^1 L^{-r})$ above by $h^0(L^{4-r})$ for all $r\geq2$.  Multiplication by the nonzero Tango map $\vartheta$ supplies an injective map in the opposite direction, since $R$ is integral.  Equality of finite dimensions proves the first asserted equality of subspaces.

Let $D=\sum D_i$ be the branch divisor of the multiroot cover.  The same logarithmic calculation gives
\[
                       \Omega_X^1\hookrightarrow\rho^*\Omega_R^1(\log D).
\]
For every $r>0$, the residue sequence and antiampleness on each $D_i$ give
\[
 H^0(R,\Omega_R^1(\log D)\otimes L^{-r})
       =H^0(R,\Omega_R^1\otimes L^{-r}).
\]
Projection formula and the first equality therefore bound
$h^0(\Omega_X^1 A^{-d})$ above by
\[
 \sum_{0\leq a_i<445}h^0(R,L^{4-d-\sum a_i})=h^0(X,A^{4-d}).
\]
The pulled-back $\vartheta$ is nonzero at the generic point, because the root cover is separable.  Multiplication by it supplies the reverse injection, proving equality as before.

Finally, if $\vartheta|_Y$ were zero, the nonzero section of $\Omega_X^1\otimes A^{-4}$ defining $\vartheta$ would factor through $\Omega_X^1\otimes A^{-(4+d)}$.  Its global sections vanish, since $4+d\geq5$.  This contradiction proves the final assertion; it does not presume that the pulled-back map is a subbundle at every ramification point.
\end{proof}

\begin{lemma}[The integral first crystalline lattice comes from the base curve]
\label{lem:OC-integral-first-lattice}
\status{Proved; independent root and reduction-descent proof, primary-source, and computation audits completed.}
Let $h:X\to C$ be the composite map in the explicit construction.  Then $h_*\OO_X=\OO_C$, its geometric generic fibre is integral, and pullback is an integral isomorphism
\[
                 H^1_{\cris}(C/W)\xrightarrow{\sim}H^1_{\cris}(X/W).
\]
For every smooth section $Y\in|A^d|$ with $d\geq2$, the induced morphism $h_Y:Y\to C$ satisfies $(h_Y)_*\OO_Y=\OO_C$, and pullback gives integral isomorphisms
\[
       H^1_{\cris}(C/W)\xrightarrow{\sim}H^1_{\cris}(Y/W),
       \qquad H^1_{\cris}(X/W)\xrightarrow{\sim}H^1_{\cris}(Y/W).
\]
\end{lemma}
\begin{proof}
First, each Tango step is flat with geometrically integral fibres, as follows from its finite locally free normalized algebra and its explicit cuspidal fibre.  Their composite $R\to C$ therefore has geometrically integral fibres.  To see integrality directly, base-change to a geometric point of $C$: the base cusp is integral, the second morphism is flat, and its geometric generic fibre is integral.  Any affine chart injects into its generic fibre algebra by flatness over the base domain, so it is reduced; all nonempty open charts meet the same integral generic fibre, giving irreducibility.  Hence $k(R)/k(C)$ is a regular field extension.

Consider the Stein factorization $X\to C'\to C$.  Because $X$ is normal, $C'$ is a smooth projective curve.  Its function field is the algebraic closure of $k(C)$ in $k(X)$.  Regularity of $k(R)/k(C)$ makes it linearly disjoint from $k(C')$; consequently
\[
                  [k(C'):k(C)]\mid[k(X):k(R)]=445^{843}.
\]
The finite cover $C'\to C$ is therefore separable.  Stein factorization gives geometrically connected fibres for $X\to C'$, so proper base change and Leray inject $H^1_{\et}(C',\mathbf Q_3)$ into $H^1_{\et}(X,\mathbf Q_3)$.  Its dimension is at most $842=2g(C)$.  Riemann--Hurwitz, including its effective different in the possibly wild case, gives
\[
                  g(C')-1\geq\deg(C'/C)\,(421-1).
\]
Thus $\deg(C'/C)=1$.  This proves $h_*\OO_X=\OO_C$.
Also $k(X)/k(C)$ is separably generated: $k(R)/k(C)$ is regular and $k(X)/k(R)$ is separable.  Since it has no new algebraic constants, it is regular, proving that the geometric generic fibre is integral.  Exact primary locators are
\cite[Theorem~37.53.4, Tag~03H0; Lemma~53.12.2, Tag~0C1D, in Section~53.12, Tag~0C1B]{StacksProject}.

We now identify the integral lattice, not just its rational span.  Both Tango steps have sections.  Their pullbacks on reduced Picard varieties are split monomorphisms; the first Betti numbers are equal, so these monomorphisms are isomorphisms.  Thus
\[
                   \Pic^0_{R,\mathrm{red}}\simeq\Pic^0_C.
\]
For the finite locally free cover $\rho$, the norm satisfies
\[
            \operatorname{Nm}_\rho\circ\rho^*=[445^{843}].
\]
The norm and this identity commute with arbitrary base change, by
\cite[Lemmas~31.18.2 and~31.18.6, Tags~0BCY and~0BD2]{StacksProject}; hence they give morphisms of Picard group schemes.  The reduced Picard varieties of $R$ and $X$ have the same dimension $421$.  Their pullback is therefore an isogeny whose kernel is contained in the torsion killed by the odd integer $445^{843}$.  It induces an isomorphism of their $2$-divisible groups.  Illusie's integral functorial identification of first crystalline cohomology with the reduced Picard realization
\cite[II, Proposition~3.11 and Remark~3.11.2]{Illusie1979DRW}
proves the first integral isomorphism in the statement.  This is where the odd covering degree, rather than rational Betti equality alone, is used.

Let now $Y$ be any smooth section of degree $d\geq2$.  It dominates $C$: an ample divisor cannot be disjoint from a positive-dimensional fibre.  Over an algebraic closure of $k(C)$, it is an effective ample Cartier divisor in the integral projective surface $X_{\overline\eta}$.  This surface is Cohen--Macaulay: it is a geometric fibre of a flat morphism from a smooth variety to a smooth curve, hence locally a Cartier complete intersection.  An ample effective divisor in such a surface is connected.  Here is the precise cohomological argument: for large $n$, Serre duality and Serre vanishing give
\[
               H^1(X_{\overline\eta},\OO(-nY))=0.
\]
Since $H^0(X_{\overline\eta},\OO)=\overline{k(C)}$, the divisor sequence gives the same constants on $nY$.  If $Y_{\overline\eta}$ were disconnected, every $nY$ would have a nontrivial idempotent global function, a contradiction.  Thus the geometric generic fibre of $Y\to C$ is connected.

We show that its finite Stein map is the identity without assuming separability in advance.  Geometric generic connectedness implies that the finite curve map $C_Y\to C$ in the Stein factorization is purely inseparable.  If it had degree $2^e>1$, it would identify, up to Frobenius twist, with the $e$th relative Frobenius of a smooth curve.  Indeed a one-variable function field over a perfect field has $[K^{1/2^e}:K]=2^e$, so there is a unique purely inseparable extension of that degree inside a fixed algebraic closure.

The pullback on $H^1(\OO_C)$ for any such map has kernel of dimension at least two.  To verify the needed bound directly on our curve, the global holomorphic exact differentials $dx$ and $x^2dx=d(x^3)$ are linearly independent: their orders at infinity are $840$ and $836$.  The exact-differential sequence (or \cite[Theorem~1 at $D=0$]{Mukai2011Kodaira})
\[
 0\longrightarrow\OO_C\longrightarrow F_*\OO_C
                    \longrightarrow B_C^1\longrightarrow0
\]
therefore gives at least two independent classes in the Frobenius kernel of $H^1(C,\OO_C)$, since Frobenius on $H^0(C,\OO_C)=k$ is surjective.  Higher Frobenius powers have at least this kernel dimension.  The same dimension statement holds after any Frobenius twist.

But $h_*\OO_X=\OO_C$ makes $H^1(C,\OO_C)\to H^1(X,\OO_X)$ injective.  A nontrivial purely inseparable Stein map for $Y$ would put at least two dimensions from this subspace in the coherent restriction kernel.  By Lemma~\ref{lem:OC-threefold-negative-degrees}, that kernel has dimension one at $d=2$ and zero for $d\geq3$.  This contradiction proves $C_Y=C$ and $(h_Y)_*\OO_Y=\OO_C$.

This equality is enough to make $\Pic_C\to\Pic_Y$ a monomorphism of fppf Picard sheaves.  Indeed it remains true after every base change over $k$, by flat base change.  If the pullback of a line bundle from $C$ is trivial, or comes from the parameter base, projection formula and pushforward recover the same assertion on $C$.  The same reasoning after a local fppf trivialization handles the sheaf quotient defining the Picard functor.  Restriction to the reduced identity components is thus a closed immersion of abelian varieties.

Finally $b_1(Y)=b_1(X)=842$.  One can use the usual \'etale weak Lefschetz statement, or obtain it directly here: $X\setminus Y$ is smooth affine of dimension three, so affine vanishing and Poincar\'e duality give $H_c^1=H_c^2=0$; localization then identifies $H^1(X,\mathbf Q_3)$ and $H^1(Y,\mathbf Q_3)$, with the same SGA~4 locators as in Lemma~\ref{lem:OC-threefold-multiroot}.  The closed immersion $\Pic_C^0\to\Pic^0_{Y,\mathrm{red}}$ is between abelian varieties of dimension $421$, hence is an isomorphism.  Illusie's integral identification yields the first isomorphism for $Y$.  Factoring that map through $X$ yields the final isomorphism of first crystalline lattices.
\end{proof}

\begin{theorem}[Exact degree range for integral injectivity in the example]
\label{thm:OC-exact-integral-injectivity-degrees}
\status{Proved; independent root and reduction-descent proof, primary-source, and computation audits completed.}
For the fixed embedded smooth threefold $(X,A)$ of
Lemma~\ref{lem:OC-threefold-multiroot} and every smooth degree-$d$ hypersurface section $Y_d$, the map
\[
            H^2_{\cris}(X/W)\longrightarrow H^2_{\cris}(Y_d/W)
\]
is injective if and only if $d\geq3$.  Its kernel contains a nonzero element killed by two for $d=1,2$.
For degree one the subgroup killed by two has dimension at least $361$, and for degree two it has dimension at least one.
The corresponding full coherent first-cohomology restriction kernels have dimensions $1203$, $1$, and $0$ in degrees one, two, and every degree at least three, respectively.
\end{theorem}
\begin{proof}
The coherent dimensions are exactly those in
Lemma~\ref{lem:OC-threefold-negative-degrees}, because each smooth ample section is connected.  At $d=1$, the functorial conjugate edge embeds a $1203$-dimensional coherent kernel into the de Rham restriction kernel.  The first term of the universal-coefficient sequence has dimension $842$.  The argument of Theorem~\ref{thm:OC-integral-injectivity-counterexample} therefore gives a subgroup of the crystalline restriction kernel killed by two, of dimension at least $1203-842=361$.

For $d=2$, the conjugate edge again puts the nonzero coherent kernel into the de Rham kernel.  Compare the natural sequences
\[
 0\longrightarrow H^1_{\cris}(-/W)/2
 \longrightarrow H^1_{\dR}(-/k)
 \longrightarrow H^2_{\cris}(-/W)[2]\longrightarrow0
\]
for $X$ and $Y_2$.  Their left map is an isomorphism by
Lemma~\ref{lem:OC-integral-first-lattice}; the snake lemma consequently identifies the de Rham kernel with the right-hand crystalline kernel.  The latter is nonzero, as asserted.

It remains to prove injectivity for every $d\geq3$.  First we show that no nonzero global closed one-form on $X$ restricts to zero on $Y=Y_d$.  The conormal sequence has no global conormal sections, since $H^0(Y,A^{-d}|_Y)=0$.  Thus a one-form restricting to zero is in
$H^0(X,\Omega_X^1(-Y))$, whose inclusion in the untwisted forms multiplies by the defining section $s_Y$ of $A^d$.
For $d\geq5$ this group is zero by
Lemma~\ref{lem:OC-threefold-differential-degrees}.
For $d=3,4$, the same lemma writes any such form uniquely as
\[
                         \omega=s_Y t\vartheta,
             \qquad t\in H^0(X,A^{4-d}).
\]
If $t\ne0$, it cannot vanish identically on $Y$, because that would give a nonzero section of $A^{4-2d}$, an antiample bundle.  Trivialize $A$ locally, write $s_Y=g$ and $t=u$, and let $\vartheta_0$ be the local one-form representing the map from $A^4$.  Modulo $g$, the equation $d\omega=0$ reads
\[
                            u\,dg\wedge\vartheta_0=0.
\]
At the generic point of the integral smooth divisor $Y$, $u$ is nonzero and $dg$ is a conormal parameter.  Therefore $\vartheta|_Y$ maps to zero in $\Omega_Y^1$.  By torsion-freeness this holds on all of $Y$, so it factors through the conormal line bundle:
\[
                         A^4|_Y\longrightarrow A^{-d}|_Y.
\]
But the space of these homomorphisms is
$H^0(Y,A^{-(d+4)}|_Y)=0$.  Hence $\vartheta|_Y$ would be zero, contradicting the final assertion of
Lemma~\ref{lem:OC-threefold-differential-degrees}.  This contradiction proves that $\omega=0$.

The plain Hodge edge now proves de Rham injectivity.  For a smooth proper connected variety, the exact sequence obtained from
$\Omega^{\geq1}\to\Omega^\bullet\to\OO$ identifies
\[
 \ker\{H^1_{\dR}\longrightarrow H^1(\OO)\}
       =H^0(\Omega^1)\cap\ker(d),
\]
and the indicated global closed forms inject into de Rham cohomology, since the differential of every global function is zero.  Any de Rham restriction-kernel class has zero coherent edge because coherent restriction is injective for $d\geq3$.  It therefore comes from a closed global form; its restriction is the zero form, because closed global forms also inject into $H^1_{\dR}(Y)$.  The preceding paragraph makes it zero.  Thus $H^1_{\dR}(X)\to H^1_{\dR}(Y)$ is injective.

By Lemma~\ref{lem:OC-integral-first-lattice}, the left map in the same universal-coefficient diagram is an isomorphism.  Its snake lemma now gives injectivity on $H^2_{\cris}[2]$.  Rational crystalline weak Lefschetz makes the full restriction kernel a finite torsion module; every nonzero such $2$-primary module contains a nonzero element killed by two.  The full kernel is consequently zero.  The exact primary rational statement is
\cite[Appendix~9, Theorem~9.1 and Corollary~9.3, pp.~63--64, arXiv version~3]{NakkajimaYobuko2022}; its proof was read in full.  This proves the asserted exact threshold.
\end{proof}

\begin{remark}[The remaining cokernel and torsion-isomorphism degree questions]
\label{rem:OC-exact-degree-scope}
\status{Proved under stated hypotheses}
The exact transition-matrix and arithmetic certificate \texttt{A\_injectivity\_degree\_checks.py} returned PASS in $0.10251374799918267$ seconds.  It verifies the actual symmetric-square decomposition, detects the invalid symmetric-dual substitution, and checks all stated finite numerical contributions; the all-degree assertions use the proofs above.

The theorem determines the full degree range for A.1(c) on this fixed example.  It also determines the full coherent first-cohomology restriction range.  It does not identify all cokernel torsion or the full degree-two torsion groups.  The exact-degree questions for A.1(a) and A.1(b), including the earlier product examples, remain distinct; an isomorphism only on subgroups killed by two would not settle them.
\end{remark}
